\section{Auxiliary results}
\label{app:aux}
\renewcommand{\theequation}{A.\arabic{equation}}
\setcounter{equation}{0}

This appendix collects the technical tools used by the proofs of
Appendix~\ref{app:proofs}: elementary probabilistic lemmas
(\ref{app:elementary}), the projection lemmas behind
Proposition~\ref{prop:projection} (\ref{app:projlem}), the deterministic
geometry of empirical-rank windows (\ref{app:rankgeom}), the local Hill
representation and the finite-sample uniform profile bound
(\ref{app:hill}), and the passage from profiles to scores, ending with a
deterministic separation lemma (\ref{app:scores}).

Throughout, the dimension is written $p=p_n$, and two auxiliary moduli
quantify the content of (E2)--(E3):
\[
 b_n(a)=\max_{j\le p}\sup_{u\in\I}\sup_{s\ge a^{-1},\,t>1}
 \frac{|\log\{\ell_j(ts,u)/\ell_j(s,u)\}|}{\log t},
 \qquad
 \omega_n(r,\tau)=\max_{j\le p}
 \sup_{u\in\I}\;
 \sup_{\substack{v\in[0,1]\\|v-u|\le r}}\;
 \sup_{w\in[\tau,1-\tau]}
 \left|\log\frac{Q_j(w^{-1},v)}{Q_j(w^{-1},u)}\right|,
\]
so that (E2) reads $b_n(3\alpha)=o(\Delta_{\min})$ and (E3) reads
$\omega_n(C_0h,\tau_n)=o(\Delta_{\min})$ with $\tau_n=(pn)^{-2}$;
$\omega_n(\cdot,\tau)$ is nondecreasing in its first argument, so (E3)
bounds $\omega_n(r,\tau_n)$ for every $r\le C_0h$.  Put
$\delta_n=(n+1)^{-1}$,
\[
 m_n^-=\bigl\lfloor\{h+\min(h,\varepsilon)\}(n+1)\bigr\rfloor-1,
 \qquad
 K_n=\lfloor\alpha m_n^-\rfloor,
 \qquad
 S_n=4n+1.
\]
All conditional distributions are the pointwise compatible versions of
Section~\ref{sec:population}, and constants are uniform over the
triangular array.

\subsection{Elementary probabilistic lemmas}\label{app:elementary}

The first lemma transfers pointwise perturbations to order statistics, as
in \citet[Appendix, Lemma~3]{gardesPodgorny2025}.

\begin{lemma}[Sorting is sup-norm nonexpansive]\label{lem:app-sorting}
Let $a_1,\ldots,a_m$ and $b_1,\ldots,b_m$ satisfy $a_i\le b_i$ for all $i$.
Then $a_{(i)}\le b_{(i)}$ for all $i$.  Consequently
$|a_i-b_i|\le\epsilon$ for all $i$ implies
$|a_{(i)}-b_{(i)}|\le\epsilon$ for all $i$.
\end{lemma}

\begin{proof}
Fix $i$.  At least $m-i+1$ indices $r$ satisfy $a_r\ge a_{(i)}$, and for
each such $r$, $b_r\ge a_r\ge a_{(i)}$; hence $b_{(i)}\ge a_{(i)}$.  For
the second claim apply the first to $(a_i-\epsilon,b_i)$ and
$(b_i-\epsilon,a_i)$.
\end{proof}

\begin{lemma}[Chernoff bound for exponential means]\label{lem:app-chernoff}
Let $L_k$ be the mean of $k$ i.i.d.\ standard exponentials.  For
$0<t\le1$,
$\Pp(|L_k-1|>t)\le2e^{-kt^2/4}$.
\end{lemma}

\begin{proof}
Markov's inequality applied to $e^{\lambda kL_k}$, optimized at
$\lambda=t/(1+t)$, gives
$\Pp(L_k>1+t)\le e^{-k\{t-\log(1+t)\}}$.  The function
$\varphi(t)=t-\log(1+t)-t^2/4$ has $\varphi(0)=0$ and
\[
 \varphi'(t)=1-\frac1{1+t}-\frac t2=\frac{t(1-t)}{2(1+t)}\ge0
 \quad\text{on }[0,1],
\]
so $t-\log(1+t)\ge t^2/4$ on $(0,1]$.  For the lower tail,
$\Pp(L_k<1-t)\le e^{-k\{-t-\log(1-t)\}}$ and
$-t-\log(1-t)\ge t^2/2\ge t^2/4$.
\end{proof}

\begin{lemma}[Threshold order statistic of uniforms]\label{lem:app-threshold}
Let $V_{(1)}<\cdots<V_{(m)}$ be the order statistics of $m$ i.i.d.\
standard uniforms, $0<\alpha\le1/3$ and
$k=\lfloor\alpha m\rfloor\ge1$.  Then
$\Pp\{V_{(k+1)}>3\alpha\}\le e^{-(k+1)/4}$.
\end{lemma}

\begin{proof}
If $3\alpha\ge1$ the probability is zero.  Otherwise
$\{V_{(k+1)}>3\alpha\}=\{N\le k\}$ with
$N\sim\mathrm{Bin}(m,3\alpha)$ and mean $\mu=3\alpha m\ge3k$.  The
multiplicative Chernoff bound
$\Pp\{N\le(1-\delta)\mu\}\le e^{-\delta^2\mu/2}$
\citep[Section~2.2]{boucheronLugosiMassart2013} with $\delta=2/3$ gives
$\Pp(N\le\mu/3)\le e^{-2\mu/9}\le e^{-2k/3}\le e^{-(k+1)/4}$, the last
step because $8k\ge3(k+1)$ for $k\ge1$.
\end{proof}

\subsection{Projection lemmas}\label{app:projlem}

Let $\mathcal G\subseteq\sigma(U)$, $T=\gamma(U)$, and
$\xi=\esssup(T\mid\mathcal G)$, the smallest $\mathcal G$-measurable random
variable dominating $T$ almost surely.

\begin{lemma}[Uniform remainder replacement]\label{lem:app-remainder}
Define $N(y)=\E\{c(U)y^{-1/T}\mid\mathcal G\}$.  Under (C1)--(C2), almost
surely,
\begin{equation}
 \left|\frac{\Pp(Y>y\mid\mathcal G)}{N(y)}-1\right|
 \le\rho(y):=\sup_u\left|\frac{L(y,u)}{c(u)}-1\right|\longrightarrow0.
 \label{eq:A1}
\end{equation}
\end{lemma}

\begin{proof}
By the tower property,
$\Pp(Y>y\mid\mathcal G)=\E\{y^{-1/T}L(y,U)\mid\mathcal G\}$.  Writing
$L(y,U)=c(U)\{1+r(y,U)\}$ with $|r(y,U)|\le\rho(y)$,
\[
 |\Pp(Y>y\mid\mathcal G)-N(y)|
 =\bigl|\E\{c(U)y^{-1/T}r(y,U)\mid\mathcal G\}\bigr|
 \le\rho(y)N(y),
\]
and $\rho(y)\to0$ is (C2).
\end{proof}

The next lemma is the point at which the argument departs from the proof
of Proposition~1 in \citet{gardesPodgorny2025}: their condition (C.1)
assumes the law of $\gamma(X)$ absolutely continuous, which identifies
$\{\gamma>t\}$ with $\{\gamma\ge t\}$ in their computation.  The tilting
argument below needs no such assumption and allows purely atomic
conditional laws of $T$ --- the relevant case for inactive coordinates
and constant-index null models.

\begin{lemma}[Tilted-measure concentration]\label{lem:app-tilt}
For $y>1$ define the random probability measure
\[
 K_y^\star(A)=
 \frac{\E\{c(U)y^{-1/T}\mathbf 1_{\{T\in A\}}\mid\mathcal G\}}{N(y)},
 \qquad A\in\mathcal B([\underline\gamma,\Gamma]).
\]
On a common probability-one event, $K_y^\star\Rightarrow\delta_\xi$ weakly
as $y\to\infty$.
\end{lemma}

\begin{proof}
Since $K_y^\star\{T>\xi+\delta\}=0$ for every $\delta>0$ by definition of
$\xi$, it suffices to prove, for fixed
$\delta\in(0,\underline\gamma)$,
\begin{equation}
 K_y^\star\{T\le\xi-\delta\}\longrightarrow0\qquad\text{almost surely.}
 \label{eq:A2}
\end{equation}
Split by the $\mathcal G$-measurable event
$G_\delta=\{\xi\ge\underline\gamma+\delta\}$.  On $G_\delta^c$,
$\xi-\delta<\underline\gamma\le T$, so $\{T\le\xi-\delta\}$ is empty and
\eqref{eq:A2} is trivial.  On $G_\delta$, set $r=\xi-\delta$,
$s=\xi-\delta/2$ and $\pi=\Pp(T>s\mid\mathcal G)$.  Then $\pi>0$ almost
surely on $G_\delta$: otherwise
$\xi'=\xi-(\delta/2)\mathbf 1_{G_\delta\cap\{\pi=0\}}$ would be a
$\mathcal G$-measurable almost-sure upper bound for $T$, strictly smaller
than $\xi$ with positive probability, contradicting minimality of the
conditional essential supremum.  Bounding numerator and denominator
separately, using only $c_-\le c\le c_+$ and monotonicity of
$g\mapsto y^{-1/g}$,
\[
 K_y^\star\{T\le r\}
 \le\frac{c_+y^{-1/r}}{c_-\pi y^{-1/s}}
 =\frac{c_+}{c_-\pi}\,y^{-(1/r-1/s)}\longrightarrow0.
\]
A rational sequence $\delta\downarrow0$ and a countable intersection give
a common probability-one event.
\end{proof}

\begin{lemma}[Continuity of the envelope]\label{lem:app-envelope-cont}
Under (C1) and (S), so that $\xi_j(u)$ equals the fibre maximum
$\max_v\gamma(u,v)$, each $\xi_j$ is continuous on $[0,1]$, with
modulus bounded by the modulus of continuity of $\gamma$.
\end{lemma}

\begin{proof}
$\gamma$ is uniformly continuous on the compact cube; if
$|u-u'|\le\delta$ then $|\gamma(u,v)-\gamma(u',v)|\le\varpi(\delta)$ for
every $v$, and taking maxima preserves the bound.
\end{proof}

\subsection{Empirical-rank geometry}\label{app:rankgeom}

The three lemmas of this subsection are purely deterministic or
rank-based; no property of the response enters.  They replace the
covariate-density conditions (H.1)--(H.2) of
\citet{gardesPodgorny2025} in the present empirical-rank setting.

\begin{lemma}[Rank displacement]\label{lem:app-dkw}
Let $\widehat U_{ij}=R_{ij}/(n+1)$ and
\[
 D_n(\eta)=\Bigl\{\max_{j\le p}\max_{i\le n}
 |U_{ij}-\widehat U_{ij}|\le\eta+\delta_n\Bigr\}.
\]
Then $\Pp\{D_n(\eta)^c\}\le2pe^{-2n\eta^2}$.
\end{lemma}

\begin{proof}
Fix $j$ and let $F_{n,j}$ be the empirical distribution function of the
$U_{ij}$, $i\le n$.  By the Dvoretzky--Kiefer--Wolfowitz inequality with
Massart's constant \citep{massart1990},
$\sup_t|F_{n,j}(t)-t|\le\eta$ outside probability
$2e^{-2n\eta^2}$.  On that event, continuity gives
$F_{n,j}(U_{(r)j})=r/n$, hence $|U_{(r)j}-r/n|\le\eta$ and
\[
 |U_{(r)j}-r/(n+1)|\le\eta+\frac{r}{n(n+1)}\le\eta+\delta_n
\]
for every rank $r$ simultaneously.  A union bound over $j$ finishes; no
union over ranks is needed.
\end{proof}

\begin{lemma}[Deterministic window counts]\label{lem:app-counts}
For every $j$, every realization and every $u\in\I$,
$M_j(u)\ge m_n^-$ and $k_j(u)\ge K_n$; moreover $M_j(u)$ is a nonrandom
function of $u$.
\end{lemma}

\begin{proof}
The ranks form a permutation of $1,\ldots,n$, so
$M_j(u)=\#\{r\in\{1,\ldots,n\}:|r/(n+1)-u|\le h\}$, a deterministic
function of $u$: the number of integers in
$[x(u),y(u)]$ with $x(u)=\max\{(u-h)(n+1),1\}$ and
$y(u)=\min\{(u+h)(n+1),n\}$.  Fix $u\in\I$ and write
$m^-=m_n^-=\lfloor\{h+\min(h,\varepsilon)\}(n+1)\rfloor-1$; note
$h+\min(h,\varepsilon)\le\min(2h,\,h+\varepsilon)$.  Four cases.

If neither endpoint is clipped, the interval has length $2h(n+1)$ and
contains at least $\lfloor2h(n+1)\rfloor-1\ge m^-$ integers.

If only the left endpoint is clipped, then $x(u)=1$ and the count is
$\lfloor y(u)\rfloor$.  Either $y(u)=(u+h)(n+1)\ge(\varepsilon+h)(n+1)$,
so the count is at least
$\lfloor(h+\varepsilon)(n+1)\rfloor\ge m^-+1$; or $y(u)=n$ and the count
is $n\ge\lfloor(h+\varepsilon)(n+1)\rfloor-1\ge m^-$, using
$h\le1-\varepsilon$, i.e.\ $(h+\varepsilon)(n+1)\le n+1$.

If only the right endpoint is clipped, then $y(u)=n$ and the count is
$n-\lceil x(u)\rceil+1\ge n-x(u)\ge n-(1-\varepsilon-h)(n+1)
=(h+\varepsilon)(n+1)-1$, hence at least
$\lceil(h+\varepsilon)(n+1)\rceil-1\ge m^-$.

If both endpoints are clipped, the count is $n$, which is at least
$m^-$ as in the second case.  In every case $M_j(u)\ge m^-$, and
$k_j(u)=\lfloor\alpha M_j(u)\rfloor\ge\lfloor\alpha m_n^-\rfloor=K_n$.
\end{proof}

\begin{lemma}[Finite-state structure]\label{lem:app-states}
For fixed $j$ and fixed ranks, $u\mapsto\mathcal W_j(u)$ takes at most
$S_n=4n+1$ distinct values on $[0,1]$.
\end{lemma}

\begin{proof}
The membership indicator of observation $i$ changes only at
$\widehat U_{ij}\pm h$: at most $2n$ breakpoints, hence at most $2n+1$
constancy intervals plus the $2n$ breakpoint states.
\end{proof}

\begin{lemma}[Local extreme count]\label{lem:app-kn}
If $nh\ge4$ and $\alpha m_n^-\ge2$, then $m_n^-\ge nh/2$ and
\[
 \frac{n\alpha h}{4}\;\le\;K_n\;\le\;2n\alpha h+1 .
\]
\end{lemma}

\begin{proof}
$m_n^-\ge\lfloor h(n+1)\rfloor-1\ge h n-2\ge nh/2$ for $nh\ge4$, and
$K_n=\lfloor\alpha m_n^-\rfloor\ge\alpha m_n^-/2\ge n\alpha h/4$ when
$\alpha m_n^-\ge2$ (since $\lfloor x\rfloor\ge x/2$ for $x\ge2$).  For the upper bound, $m_n^-\le2h(n+1)-1\le2hn+2h-1\le2hn+1$, so
$K_n\le\alpha m_n^-\le2n\alpha h+\alpha\le2n\alpha h+1$.
\end{proof}

The rate conditions \eqref{eq:rate-cond} force $nh\to\infty$ and
$K_n\to\infty$, so Lemma~\ref{lem:app-kn} applies for $n$ large; it is
the only conversion between $K_n$ and $n\alpha h$ used in
Appendix~\ref{app:proofs}.

\subsection{Local Hill representation and uniform profile control}
\label{app:hill}

The scheme parallels the proof of Theorem~3 in
\citet{gardesPodgorny2025} --- oscillation control,
exponential-spacings representation, slowly varying remainder --- with
the covering net over their compact matrix class replaced by exact union
bounds over $p$ coordinates and $S_n$ window states.

\begin{lemma}[Conditional PIT and R\'enyi representation]\label{lem:app-pit}
Fix $j$ and condition on the column $\bm U_j=(U_{1j},\ldots,U_{nj})$.
Under (E1):
\begin{enumerate}
\item $V_{ij}=1-F_{Y|j}(Y_i\mid U_{ij})$, $i\le n$, are i.i.d.\ standard
uniform, and $Y_i=Q_j(1/V_{ij},U_{ij})$ almost surely;
\item for any index set $\mathcal W$ of size $m$, deterministic given
$\bm U_j$, with $V_{(1)}<\cdots<V_{(m)}$ the ordered PIT values of its
members and $k=\lfloor\alpha m\rfloor$,
\[
 L_k=\frac1k\sum_{r=1}^k\log\frac{V_{(k+1)}}{V_{(r)}}
\]
is distributed as the mean of $k$ i.i.d.\ standard exponentials,
independent of $V_{(k+1)}$.
\end{enumerate}
\end{lemma}

\begin{proof}
(i) Given $\bm U_j$, the pairs $(U_{ij},Y_i)$ are independent in $i$ and
$Y_i$ has continuous conditional distribution function
$F_{Y|j}(\cdot\mid U_{ij})$, so $V_{ij}$ is standard uniform.  For the
representation, $Q_j(1/V_{ij},U_{ij})=F_{Y|j}^{-1}(F_{Y|j}(Y_i\mid
U_{ij})\mid U_{ij})\le Y_i$ always, with strict inequality only when
$Y_i$ lies in the interior of an interval on which
$F_{Y|j}(\cdot\mid U_{ij})$ is constant; each such interval has
conditional probability zero and there are countably many, so the
representation holds almost surely.  Strict monotonicity is not needed.
Conditioning on the whole column, rather than on window counts as in
\citet[Appendix, Lemma~4]{gardesPodgorny2025}, makes the window sets
deterministic (Lemma~\ref{lem:app-counts}); only coordinate $j$ is
conditioned on, and no independence across coordinates is used.

(ii) With $W_i=-\log V_i$ standard exponentials and order statistics
$W_{(1)}<\cdots<W_{(m)}$, $W_{(m-r+1)}=-\log V_{(r)}$, Abel summation
gives
\[
 kL_k=\sum_{r=1}^k\{W_{(m-r+1)}-W_{(m-k)}\}
 =\sum_{i=1}^k i\,\{W_{(m-i+1)}-W_{(m-i)}\},
\]
and by the R\'enyi representation of exponential order statistics
\citep{renyi1953} the normalized top spacings
$i\{W_{(m-i+1)}-W_{(m-i)}\}$, $i\le k$, are i.i.d.\ standard exponential
and independent of $W_{(m-k)}$, hence of $V_{(k+1)}$.
\end{proof}

\begin{lemma}[Frozen-location Hill expansion]\label{lem:app-hill}
Fix $j$, condition on $\bm U_j$, and fix a window state $\mathcal W$ of
size $m$ realized at $u\in\I$, with $k=\lfloor\alpha m\rfloor$ and Hill
statistic $H$.  On the event $D_n(\eta)$ intersected with
$\{V_{ij}\in[\tau,1-\tau]\ \forall i\}$ and $\{V_{(k+1)}\le3\alpha\}$,
\begin{equation}
 |H-\xi_j(u)|
 \le2\omega_n(h+\eta+\delta_n,\tau)+\xi_j(u)|L_k-1|+b_n(3\alpha)L_k.
 \label{eq:A3app}
\end{equation}
\end{lemma}

\begin{proof}
\emph{Spatial freezing.}  Each $i\in\mathcal W$ has
$|\widehat U_{ij}-u|\le h$, hence $|U_{ij}-u|\le r_n:=h+\eta+\delta_n$ on
$D_n(\eta)$.  With $V_{ij}\in[\tau,1-\tau]$, the definition of
$\omega_n$ gives
$|\log Q_j(1/V_{ij},U_{ij})-\log Q_j(1/V_{ij},u)|\le\omega_n(r_n,\tau)$
for every member.  Let $\check H$ be the Hill statistic of the frozen
responses $\check Y_i=Q_j(1/V_{ij},u)$.  By Lemma~\ref{lem:app-sorting},
ordered log-responses differ by at most $\omega_n(r_n,\tau)$, and the Hill
statistic is an average of differences of two of them, so
$|H-\check H|\le2\omega_n(r_n,\tau)$.

\emph{Exact decomposition.}  Since $s\mapsto Q_j(s,u)$ is nondecreasing,
$\check Y_{(m-r+1)}=Q_j(1/V_{(r)},u)$, and
$\log Q_j(s,u)=\xi_j(u)\log s+\log\ell_j(s,u)$ yields
\[
 \check H=\xi_j(u)L_k+R,
 \qquad
 R=\frac1k\sum_{r=1}^k
 \log\frac{\ell_j(1/V_{(r)},u)}{\ell_j(1/V_{(k+1)},u)}.
\]

\emph{Remainder.}  On $\{V_{(k+1)}\le3\alpha\}$, each ratio is
$\ell_j(ts,u)/\ell_j(s,u)$ with $s=1/V_{(k+1)}\ge(3\alpha)^{-1}$ and
$t=V_{(k+1)}/V_{(r)}>1$, so
$|\log(\cdot)|\le b_n(3\alpha)\log(V_{(k+1)}/V_{(r)})$; averaging gives
$|R|\le b_n(3\alpha)L_k$.
\end{proof}

The decomposition behind \eqref{eq:A3app} is the shape of the whole
argument:
\[
 \widehat\xi_j(u)-\xi_j(u)
 =\text{rank/spatial error}
 +\text{exponential fluctuation}
 +\text{tail remainder}.
\]

\begin{proposition}[Finite-sample uniform profile bound]
\label{thm:app-rank-profile}
Assume (C1)--(C2), (S), the structural part of (E1) (continuity of
every $F_{Y|j}(\cdot\mid u)$, hence the quantile representation), that
$b_n(3\alpha)$ and $\omega_n(\cdot,\cdot)$ are finite,
$0<h\le1-\varepsilon$, $0<\alpha\le1/3$, and $K_n\ge2$.  For every
$\eta>0$, $0<t\le1$, $0<\tau<1/2$,
\begin{align}
&\Pp\left\{\max_{j\le p}\sup_{u\in\I}
 |\widehat\xi_j(u)-\xi_j(u)|
 >2\omega_n(h+\eta+\delta_n,\tau)+\Gamma t+(1+t)b_n(3\alpha)\right\}
 \nonumber\\
&\qquad\le
 \underbrace{2pe^{-2n\eta^2}}_{\text{ranks}}
 +\underbrace{2pn\tau}_{\text{PIT truncation}}
 +\underbrace{2pS_ne^{-K_nt^2/4}}_{\text{spacings}}
 +\underbrace{pS_ne^{-(K_n+1)/4}}_{\text{threshold}}.
 \label{eq:A4app}
\end{align}
\end{proposition}

\begin{proof}
Call a window state \emph{admissible} for coordinate $j$ if it equals
$\mathcal W_j(u)$ for some $u\in\I$; by
Lemmas~\ref{lem:app-counts} and~\ref{lem:app-states} there are at most
$S_n$ admissible states per coordinate and each has
$k\ge K_n\ge2$, so $L_k$ and $V_{(k+1)}$ are well defined on every one of
them.  (States realized only outside $\I$ may have smaller windows; they
are not needed.)  Define the exceptional events
$T_1=D_n(\eta)^c$;
$T_2=\{\exists\,i,j:V_{ij}\notin[\tau,1-\tau]\}$;
$T_3=\{\exists\,j,\text{ admissible state }\mathcal W:|L_k-1|>t\}$;
$T_4=\{\exists\,j,\text{ admissible state }\mathcal W:
V_{(k+1)}>3\alpha\}$.

Outside $T_1\cup T_2\cup T_3\cup T_4$, fix $j$ and $u\in\I$: the realized
window is an admissible state, and
Lemma~\ref{lem:app-hill} with $\xi_j\le\Gamma$, $|L_k-1|\le t$,
$L_k\le1+t$ gives the displayed deviation bound; the modulus
$\omega_n(h+\eta+\delta_n,\tau)$ covers every selected observation
because its definition allows the displaced point $v=U_{ij}$ to range
over all of $[0,1]$.

Probabilities: $\Pp(T_1)\le2pe^{-2n\eta^2}$ is
Lemma~\ref{lem:app-dkw}; $\Pp(T_2)\le2pn\tau$ since each $V_{ij}$ is
uniform.  For $T_3$, $T_4$, condition on $\bm U_j$; the admissible states
are then deterministic with $k\ge K_n$,
Lemma~\ref{lem:app-pit} applies to each, and
Lemmas~\ref{lem:app-chernoff} and~\ref{lem:app-threshold} give conditional
bounds $2e^{-kt^2/4}\le2e^{-K_nt^2/4}$ and
$e^{-(k+1)/4}\le e^{-(K_n+1)/4}$, both decreasing in $k$.  A union bound
over the at most $pS_n$ admissible states and
integration over $\bm U_j$ complete the proof; overlapping states and
dependent coordinates cost nothing beyond the union bound.
\end{proof}

\subsection{From profiles to scores}\label{app:scores}

Recall the score grid $G=\{r\delta_n:r=1,\ldots,n\}\cap\I$, of size
$L_n=|G|$, and that $K_n\ge2$ makes every grid point usable, so
$\widehat\Psi_j=L_n^{-1}\sum_{u\in G}\widehat\xi_j(u)$.

\begin{lemma}[Quadrature]\label{lem:app-quadrature}
If $L_n\ge1$, then for every $j$,
\begin{equation}
 \Bigl|L_n^{-1}\sum_{u\in G}\xi_j(u)-\Psi_j\Bigr|
 \le q_n:=\omega_\xi(\delta_n)+\frac{4\Gamma\delta_n}{1-2\varepsilon}.
 \label{eq:A5app}
\end{equation}
\end{lemma}

\begin{proof}
Write the grid points $u_1<\cdots<u_{L_n}$, consecutive multiples of
$\delta_n$.  If $L_n=1$, then $\I$ contains a single multiple of
$\delta_n$, which forces $|\I|<2\delta_n$; the left side is at most
$\Gamma$, while $4\Gamma\delta_n/|\I|>2\Gamma$, so \eqref{eq:A5app}
holds trivially.  Assume $L_n\ge2$ and partition $\I$
exactly: $C_1=[\varepsilon,u_1+\delta_n/2]$,
$C_{L_n}=[u_{L_n}-\delta_n/2,1-\varepsilon]$, and
$C_r=[u_r-\delta_n/2,u_r+\delta_n/2]$ otherwise.  Interior cells have
length $\delta_n$, boundary cells length in $[\delta_n/2,3\delta_n/2)$,
and every point of $C_r$ is within $\delta_n$ of $u_r$.  Then
\[
 \Bigl|\frac1{L_n}\sum_r\xi_j(u_r)-\frac1{|\I|}\int_{\I}\xi_j\Bigr|
 \le\frac1{|\I|}\sum_r\int_{C_r}|\xi_j-\xi_j(u_r)|
 +\Gamma\sum_r\Bigl|\frac{\lambda(C_r)}{|\I|}-\frac1{L_n}\Bigr|.
\]
The first term is at most $\omega_\xi(\delta_n)$.  For the second, both
weight vectors sum to one;
$|\I|=\sum_r\lambda(C_r)\in[(L_n-1)\delta_n,(L_n+1)\delta_n]$ gives
$||\I|/L_n-\delta_n|\le\delta_n/L_n$, so interior cells contribute at most
$\delta_n$ in total and the two boundary cells at most
$\delta_n/2+\delta_n/L_n\le\delta_n$ each; the sum is at most $3\delta_n$,
and the second term is at most $3\Gamma\delta_n/|\I|\le
4\Gamma\delta_n/|\I|$.
\end{proof}

\begin{proposition}[Finite-sample score bound]\label{thm:app-score}
Under the conditions of Proposition~\ref{thm:app-rank-profile} (with in
addition $\omega_\xi(\delta_n)<\infty$) and $L_n\ge1$, for every
$\eta>0$, $0<t\le1$, $0<\tau<1/2$,
\begin{align}
 \Pp\Bigl\{\max_{j\le p}|\widehat\Psi_j-\Psi_j|
 >B_n(\eta,t,\tau)+q_n\Bigr\}
 \le
 2pe^{-2n\eta^2}+2pn\tau+2pS_ne^{-K_nt^2/4}+pS_ne^{-(K_n+1)/4},
 \label{eq:A6app}
\end{align}
where
$B_n(\eta,t,\tau)=2\omega_n(h+\eta+\delta_n,\tau)+\Gamma t
+(1+t)b_n(3\alpha)$.
\end{proposition}

\begin{proof}
Pathwise,
\[
 |\widehat\Psi_j-\Psi_j|
 \le\sup_{u\in\I}|\widehat\xi_j(u)-\xi_j(u)|
 +\Bigl|L_n^{-1}\sum_{u\in G}\xi_j(u)-\Psi_j\Bigr|,
\]
and Lemma~\ref{lem:app-quadrature} bounds the second term by $q_n$ on
every sample path; reusing the same observations for estimation and
averaging therefore costs no extra probability.
Proposition~\ref{thm:app-rank-profile} gives \eqref{eq:A6app}.
\end{proof}

\begin{lemma}[Deterministic score separation]\label{lem:app-separation}
Suppose $\Delta_{\min}>0$ and, on some event,
\[
 \max_{j\le p}\,|\widehat\Psi_j-\Psi_j|<\frac{\Delta_{\min}}{2}.
\]
Then, on that event, every active coordinate has a strictly smaller
score than every inactive coordinate; consequently
$\A\subseteq\widehat\A_d$ for every $d\ge s$, and
$\widehat\A_s=\A$.
\end{lemma}

\begin{proof}
Inactive coordinates have $\Psi_j=\gmax$
(Proposition~\ref{prop:geometry}, proved independently of this lemma
in Appendix~B), hence
$\widehat\Psi_j>\gmax-\Delta_{\min}/2$; active coordinates have
$\Psi_j\le\gmax-\Delta_{\min}$, hence
$\widehat\Psi_j<\gmax-\Delta_{\min}/2$.  Every active score therefore
precedes every inactive score in the increasing order, and the screening
and recovery statements follow from the definition of
$\widehat\A_d$.
\end{proof}
